% AY2012 材料力学 〔II〕（転記）
% 出典: 04_材料力学/original/04_材料力学/2012/2012za.pdf
% 要約: 上端を固定された板厚 t 一定の台形板（上端幅 b2, 下端幅 b1）の
%       引張り変形。座標 x の原点は棒の下端 B 点。自重は無視。
%       下端 B に引張荷重 P。A(x), N(x)/σ(x), 全体の伸び λ を問う。

\section*{〔II〕}

図に示す上端を固定された板厚一定の台形板（上端と下端の幅 $b_2$, $b_1$）の
引張り変形に関する問いに答えよ. 自重は無視せよ. 座標 $x$ の原点を棒の下端 B 点にとるものとする.

\begin{center}
\begin{tikzpicture}[>=Latex]
  % 台形板（上が広い b2, 下が狭い b1）。原点 B は下端。
  \def\Hh{5}      % 長さ L
  \def\Wt{2.4}    % 上端半幅
  \def\Wb{0.5}    % 下端半幅
  \draw[thick] (-\Wt,\Hh) -- (\Wt,\Hh) -- (\Wb,0) -- (-\Wb,0) -- cycle;
  % 上端固定（ハッチング）
  \draw[pattern=north east lines] (-\Wt-0.2,\Hh) rectangle (\Wt+0.2,\Hh+0.25);
  % 上端幅 b2
  \draw[<->] (-\Wt,\Hh+0.55) -- (\Wt,\Hh+0.55) node[midway,fill=white]{幅 $b_2$};
  % 下端幅 b1
  \draw[<->] (-\Wb,-0.5) -- (\Wb,-0.5) node[midway,fill=white,font=\scriptsize]{幅 $b_1$};
  % 板厚 t
  \node at (0,\Hh*0.62) {板厚 $t$};
  % 座標 x における幅 b(x)
  \draw[<->] (-1.25,1.8) -- (1.25,1.8) node[midway,fill=white,font=\scriptsize]{幅 $b(x)$};
  % 長さ L 寸法
  \draw[<->] (-\Wt-0.9,0) -- (-\Wt-0.9,\Hh) node[midway,fill=white]{$L$};
  % 座標 x（下端 B から上向き）
  \draw[->] (\Wt+0.6,0) -- (\Wt+0.6,2.2) node[above]{$x$};
  % 節点ラベル
  \node at (\Wt+0.35,\Hh) {A};
  \node at (-0.7,-0.15) {B};
  % 荷重 P
  \draw[->,very thick] (0,-0.55) -- (0,-1.8) node[midway,right]{$P$};
\end{tikzpicture}
\end{center}

\begin{enumerate}[label=(\arabic*)]
  \item 座標 $x$ における断面積 $A(x)$ を求めよ.

  \item 荷重 $P$ により座標 $x$ に生じる内力 $N(x)$ および応力 $\sigma(x)$ を求めよ.

  \item 棒全体の変位（伸び）$\lambda$ を求めよ.
\end{enumerate}
